\section{Introduction}
\label{sec:introduction}

How can a learning system construct, compose, and improve its own internal
computational mechanisms? Neural architectures already learn representations,
select information, implement memory, and organize reusable computation.
Attention, for example, learns input-dependent interactions through shared
projections and weighted aggregation~\cite{vaswani2017attention}. The question
is how computational capabilities can accumulate under continued learning:
which internal operations are acquired, how later tasks reuse them, and what
must be preserved when the parameters implementing those operations change.
An adequate answer should expose the executed computation and its learning
history, together with the resources required to retain it.

This problem has two coupled scales. At the implementation scale, a shared
parameter update changes representations, operations, and their invocation.
At the semantic scale, a capability is a complete input--computation--output
relation. Preserving an isolated operator is insufficient if its old inputs
are subsequently encoded differently or routed elsewhere. Conversely,
preserving all implementation details can prevent useful adaptation or require
separate capacity for each task. Dynamically expandable networks illustrate
that selective sharing, structural growth, and protection can be organized
jointly~\cite{yoon2018lifelong}; their existence motivates a precise comparison
of learning organizations, rather than a claim that conventional neural
networks cannot support these functions.

We investigate \emph{Prometheus}, a candidate organization in which executable
operations are represented by monodromy actions of learned loops. Continuation
around a regular loop transports an internal state to another state, and
concatenating loops composes their actions in execution order. Different
geometric implementations can realize the same operation. A certified
deformation within a loop class therefore supplies shared implementation
freedom while preserving its endpoint semantics. Acquiring a different
operation requires a different learning action: a finite controller update can
move an unadmitted loop between regular regions separated by a discriminant.
The intermediate singular path is not executed.

The learning interface couples this class-changing process to neural encoding,
readout, and operation calls. Terminal task feedback trains those bindings
through the currently executed operations. For class acquisition, a smoothed
objective assigns credit to measured loss changes across loop-class
boundaries. The resulting update combines ordinary neural derivatives with a
boundary contribution. Once a capability is admitted, geometric protection,
closure of its native input--operation relation, and a finite-length
soft-execution agreement condition address distinct retention obligations.
This organization makes the relationship between operation formation and
protected shared learning explicit and testable.

The central hypothesis is that this separation of executable semantics from
its mutable implementation can improve capability acquisition and retention
at a useful resource tradeoff. Homotopy invariance alone does not establish
that benefit. The comparison must give conventional learners the same task
information, feedback, execution observations, and usable certificates. It
must also charge operation discovery, failed proposals, state growth,
verification, storage, and execution alongside neural parameters and gradient
updates. Parameter sharing, ordinary behavioral constraints, and compilation
remain separate possible explanations of an observed gain.

We develop the mathematical learning interface and an inspectable Block
proposal around this hypothesis. An independent advantage remains
unestablished: the fully integrated architecture must be fixed and evaluated
across continual acquisition, demanding compositions, broader computational
domains, real inputs, and competitive learning systems. These requirements
connect the proposed mechanism to the original question of self-organizing
computation while keeping its scientific claims conditional on executable
evidence.
